\chapter{序論}

幾何学の公理は、かつては思考の法則とみなされ、聡明な精神であれば否定することも吟味することもできないものと考えられていた。われわれが慣れ親しんできた公理が経験とよく一致することが見いだされたばかりでなく、それらのどれかが成り立たないという仮定のもとでは推論することすらできない、と信じられていた。しかし今では、ユークリッドの公理と全体として、あるいは一部分において矛盾する公理の組を採用しても、ユークリッドの幾何学と同じように無矛盾な幾何学を築き上げることができる、ということが示されている。

ここでは、最も重要な二つの非ユークリッド幾何学を示す。\footnote{「歴史的注」（本書末）を見よ。}これらでは、平行線に関するもの以外の公理と定義は、ユークリッドのとおりに採る。平行線についての公理を除くと、三つの仮説に至る。その一つはユークリッドの幾何学を確立し、他の二つはそれぞれ、興味深くかつ有用な一連の命題を与える。実際、宇宙のうち限られた部分しか調べることができない以上、ユークリッドの体系が、これから述べる二つの非ユークリッド幾何学のいずれよりも真であることを証明することはできない。

配列は次のようにする。まず三つの幾何学に共通する命題を証明し、つぎにユークリッドと異なる二つの幾何学のそれぞれについて、一連の命題と三角法の公式を導き、最後に解析的方法によって、それらの最も顕著な性質のいくつかを導く。

幾何学の基礎を直接に調べることも、この研究で意識的にせよ無意識的にせよなされた仮定をすべて指摘することも、ここでは意図しない。これらの幾何学に共通なものには手をつけずにおくことで、それらの相違点だけに注意を集中することができる。抽象的な理論だけによるよりも、具体的な叙述によるほうが、幾何学の本性についてより多くを学べるかもしれない。

そこで、幾何学の用語の大部分は、教科書に与えられている定義をくり返すことなく用い、これらの用語で定義される図形が存在するものと仮定する。特に、次を仮定する。
\smallskip

\begin{enumerate}
\item[I.] \emph{任意の二点によって定まる直線が存在すること、および二点間の最短の道は直線であること。}

\item[II.] \emph{一直線上にない任意の三点によって定まる平面が存在すること、および平面上の任意の二点を結ぶ直線は完全にその平面に含まれること。}

\item[III.] \emph{幾何学的図形は、形も大きさも変えずに動かすことができること。}

\item[IV.] \emph{点が直線上をある位置から別の位置へ動くとき、その間のすべての点を通ること、また、たとえば角や線分の長さのような幾何学的な量がある値から別の値まで変化するとき、その間のすべての値をとること。}
\end{enumerate}

いくつかの命題では、細部を普通の教科書から補えるものについて、証明を省くか、証明の方法を示すだけにとどめる。
